% !TeX program = pdflatex
% ─────────────────────────────────────────────────────────────────────────────
%  corpo_universal.tex — CORPO UNIVERSAL (raiz)
%
%  Objecto de nível superior: o corpo das transformações entre construções.
%  Não é supercorpo. Não entra no bestiário. INGEST ⊂ U (realização).
%  Compila: pdflatex corpo_universal.tex
%
% ─────────────────────────────────────────────────────────────────────────────
\documentclass[11pt,a4paper]{article}
\usepackage[utf8]{inputenc}\usepackage[T1]{fontenc}\usepackage[brazil]{babel}
\usepackage{amsmath,amssymb,amsthm}
\usepackage[margin=2.4cm]{geometry}
\usepackage{booktabs}\usepackage{array}\usepackage{tabularx}\usepackage{xcolor}
\usepackage[htt]{hyphenat}
\usepackage[hidelinks,breaklinks]{hyperref}

\sloppy
\emergencystretch=2em

\definecolor{cinza}{gray}{0.35}
\newcommand{\code}[1]{\texttt{\small\codecolon#1:\endcolon}}
\def\codecolon#1:#2\endcolon{\codehyphen#1-\endhyphen\ifx\relax#2\relax\else:\allowbreak\codecolon#2\endcolon\fi}
\def\codehyphen#1-#2\endhyphen{#1\ifx\relax#2\relax\else-\allowbreak\codehyphen#2\endhyphen\fi}
\newcommand{\caixa}[1]{%
  \par\smallskip\noindent
  \begin{center}
  \fbox{\begin{minipage}{0.92\linewidth}\centering\smallskip #1\smallskip\end{minipage}}
  \end{center}\smallskip}
\newtheorem{teorema}{Teorema}
\newtheorem{definicao}[teorema]{Definição}
\newtheorem{corolario}[teorema]{Corolário}
\newtheorem{proposicao}[teorema]{Proposição}
\newtheorem{observacao}[teorema]{Observação}
\renewcommand{\proofname}{Demonstração}

% ─────────────────────────────────────────────────────────────────────────────
%  papers/gkcapa.tex — a capa do Reino, mínima, para os papers em `article`.
%
%  O estilo.tex completo é do `report` (títulos de capítulo, skak, …). Aqui fica
%  só o que a capa precisa, para o paper continuar a compilar sozinho com
%  pdflatex. O tradutor próprio (tests/tex) carrega o estilo.tex da raiz / o
%  symlink papers/estilo.tex e avalia o mesmo `\gkcapa`.
% ─────────────────────────────────────────────────────────────────────────────
\definecolor{ouro}{HTML}{B8912F}
\definecolor{ouroclaro}{HTML}{F3E9CE}
\definecolor{tinta}{HTML}{1E1B16}
\definecolor{regua}{HTML}{6E6858}
\providecommand{\gknota}{\fontsize{7.62}{11.05}\selectfont}
\providecommand{\gktexto}{\fontsize{10.50}{15.22}\selectfont}
\providecommand{\gksub}{\fontsize{12.33}{17.87}\selectfont}
\providecommand{\gksec}{\fontsize{14.47}{20.98}\selectfont}
\providecommand{\gkcap}{\fontsize{16.99}{24.63}\selectfont}
\providecommand{\gktit}{\fontsize{23.42}{33.95}\selectfont}
\providecommand{\Prj}{\mathbb{P}}
\providecommand{\Trf}{\mathcal{T}}
\providecommand{\gkcapa}[3]{%
\title{\vspace{-0.6cm}
{\color{ouro}\rule{\textwidth}{1.4pt}}\\[6mm]
\textbf{\gktit\color{tinta} Reino Dourado}\\[3mm]{\gksec\itshape\color{regua} Golden Kingdom}\\[8mm]
{\gkcap\scshape\color{ouro} #1}\\[5mm]
{\gksub\color{regua} #2}\\[2mm]
{\gktexto\color{regua}\itshape #3}\\[7mm]
{\color{ouro}\rule{\textwidth}{0.6pt}}\\[9mm]{\gknota\color{regua}\begin{minipage}{\textwidth}\setlength{\parindent}{0pt}%
\textbf{Aviso de Isenção Algorítmica:} O universo, as leis e as entidades contidas nestas páginas ---
incluindo felinos matriciais, aves de rapina algébricas e amuletos de controle de fluxo --- são
construções de pura ficção formal. Esta obra não descreve a realidade física, não serve como
engenharia reversa do cosmos e não constitui doutrina religiosa. Qualquer semelhança com bugs reais do seu
sistema operacional é mera coincidência (ou falta de reza). Prossiga por sua própria conta e
risco.\end{minipage}}}
\author{}\date{}}


\begin{document}
\gkcapa{O Corpo Universal}
       {espaço, dual, base das oito leis}
       {consome a Física; não inventa estrutura; o catálogo lê as coordenadas}
\maketitle

\begin{abstract}\noindent
O \textbf{Corpo Universal} \(\mathcal{U}\) é o intermediário operacional
entre a Física e o catálogo. Não é supercorpo. Não é o 14º lugar.
A estrutura não se protocola: \textbf{já está construída}.
\[
\boxed{\;\mathcal{U}=(X,X^{*},\mathsf{Mor},\mathsf{Aut})\;}
\]
\(X\) é o espaço vectorial de largura oito, com dual \(X^{*}\).
A base ortonormal \(\{e_k=2^{k}\}_{k=0}^{7}\) \textbf{é} as oito leis
(\code{fis:thm:base}; \code{papers/aranha.c}, \code{papers/aranha.asm}).
\(\mathsf{Mor}\) relaciona construções.
\(\mathsf{Aut}\) processa um vector nesse espaço.
\[
\boxed{\;\text{Física}\;\longleftarrow\;\mathcal{U}\;\longrightarrow\;\text{Catálogo}\;}
\]
Não se promove evidência a estrutura. Não há Lei~8.
\end{abstract}

\tableofcontents
\bigskip

%═══════════════════════════════════════════════════════════════════════════
\section{O que isto é, e o que não é}\label{univ:sec:estatuto}

\begin{definicao}[Corpo Universal]\label{univ:def:U}
O \textbf{Corpo Universal} \(\mathcal{U}\) consome o espaço que a Física
já construiu --- Def.~\code{fis:def:doisandares}, Def.~\code{fis:def:dual},
Teor.~\code{fis:thm:base} --- e opera nele. Não o inventa. Não contém
os treze lugares.
\[
\boxed{\;\mathcal{U}=(X,X^{*},\mathsf{Mor},\mathsf{Aut})\;}
\]
Em baixo, \(X\) é espaço vectorial sobre o binário \(B\), dimensão \(8\),
e \(X^{*}\) o dual pelo emparelhamento que torna \(\{e_k\}\) ortonormal
(Def.~\ref{univ:def:dualU}). Em cima, o estado de \(\mathsf{Aut}\) é um
vector \(G\) do módulo das funções sobre \(X\)
(Def.~\code{fis:def:doisandares}).
\(\mathsf{Mor}=\{\mathrm{Hom},\mathrm{Iso},\mathrm{Diff},\operatorname{Isom}_{\mu},\mathbf{Duo}\}\)
é preservação.
\[
\boxed{\;\text{Física}\;\longleftarrow\;\mathcal{U}\;\longrightarrow\;\text{Catálogo}\;}
\]
Não entra no bestiário. Não é uma categoria universal.
\end{definicao}

\noindent \(\mathsf{Mor}\) pergunta que contrato \(f\) preserva;
\(\mathsf{Aut}\) pergunta o estado e a transição.
\caixa{%
$\begin{array}{rcl}
\text{categoria}
& = &
\text{como as construções se relacionam}\\[3pt]
\text{autómato}
& = &
\text{como uma construção é processada}\\[3pt]
\text{Física}
& = &
\text{o que efectivamente foi construído}\\[3pt]
\text{catálogo}
& = &
\text{o registo do resultado}
\end{array}$}
\caixa{\(\mathcal{U}\neq\text{catálogo}\neq\operatorname{INGEST}\)}
Um capítulo \textbf{não} vira corpo por ser ingerido, nem por admitir
um mapa. O percurso Álgebra\(\to\cdots\to\)Redes é execução do autómato
sobre uma sequência de entradas, não composição de 2-morfismos.

\begin{observacao}[não contaminar a taxonomia]\label{univ:obs:taxonomia}
Declarar \(\mathcal{U}\) objecto de nível superior \textbf{não} implica
que todo capítulo seja um corpo. A cadeia
\(\mathbb{N}\to\mathbb{Z}\to\mathbb{Q}\to\mathbb{R}\) fecha por opostos,
quocientes e completação, com \(\mathcal{D}=\texttt{N/A}\): isso
\emph{satisfaz} o contrato de ingestão (o slot responde) e \emph{não}
realiza duomorfismo. A forma \(E^{2}-B^{2}\) da Parte EM \textbf{não} é
corpo (\code{fis:def:em-forma}) com a interface fechada.
\end{observacao}

%═══════════════════════════════════════════════════════════════════════════
\section{O espaço, o dual, e as oito leis}\label{univ:sec:espaco}

\begin{definicao}[o espaço]\label{univ:def:espaco}
A Física já o tem (Def.~\code{fis:def:doisandares}). \(\mathcal{U}\)
\textbf{consome}:
\[
\boxed{\;
X\text{ espaço vectorial sobre }B,\quad
\dim X=8,\quad
e_k=2^{k}\ (0\le k<8).
\;}
\]
As direções de aresta do hipercubo \emph{são} esta base. Não se escolhe
armazenamento: o byte \emph{é} a escrita nas oito coordenadas
(\code{fis:thm:base}(2): \(\langle b,e_k\rangle=(b\gg k)\wedge 1\)).
\end{definicao}

\begin{definicao}[o dual]\label{univ:def:dualU}
O dual de \(X\) é o emparelhamento da Def.~\code{fis:def:dual}:
\[
\langle a,b\rangle=\operatorname{paridade}(a\wedge b),
\qquad
X^{*}\simeq X
\text{ via }\{e_k\}.
\]
Em cima, o módulo \(V\) das funções tem dual \(V^{*}\)
(Def.~\code{fis:def:dual}) e caracteres
\(\chi_k(j)=(-1)^{\langle k,j\rangle}\). Três objectos, e não se misturam:
\[
\boxed{\;X^{*}\neq V^{*}\neq\mathbf{Duo}.\;}
\]
O dual linear não é a troca \((\oplus,\otimes)\mapsto(\otimes,\oplus)\)
(Obs.~\ref{univ:obs:D}).
\end{definicao}

\begin{teorema}[a base é ortonormal, e é as oito leis]\label{univ:thm:base}
\[
\boxed{\;\langle e_i,e_j\rangle=\delta_{ij}\;}
\]
--- Teor.~\code{fis:thm:base}, Gram \(=\mathrm{Id}\). Donde a coordenada
é ler o bit. No caso do byte (\(m=8\)), as oito máscaras \emph{são} os
\(e_k\), e identificam-se com as oito leis da torre
(\code{fis:obs:octoniao-interface}: leis \(0\)--\(7\); não há Lei~8):
\begin{center}\small
\begin{tabular}{@{}cll@{}}
\toprule
\(k\) & lei & o que a Física já nomeia \\
\midrule
0 & Lei~0 & \((+1)\oplus(-1)\); \(0^{\dagger}=\infty\) \\
1 & Lei~1 & \(\nu\circ\nu=\mathrm{id}\) \\
2 & Lei~2 & \(K^{**}=K\) \\
3 & Lei~3 & \(\{-1,0,+1\}\) \\
4 & Lei~4 & \(T+T^{*}\); \(\lvert\det\rvert=1\) \\
5 & Lei~5 & \(x^{2}=-1\) (bit \(i\); \textbf{não} a estrela) \\
6 & Lei~6 & interface \(\oplus=\otimes\) no suporte em que isso está escrito \\
7 & Lei~7 & \(H\times H^{*}\); dois tecidos ligam sem fundir \\
\bottomrule
\end{tabular}
\end{center}
A Lei~4 \emph{é} a métrica canónica de \(\mathcal{U}\)
(Def.~\ref{univ:def:volume}). A Lei~5 \textbf{não} é
\(\operatorname{Star}\) do elemento (\(x^{2}=x+1\)).
O dual explícito conhecido no medidor é só Lei~0: \(0^{\dagger}=\infty\).
Os demais pares \textbf{não} se inventam (\code{papers/aranha.c}).
\end{teorema}
\begin{proof}
É o Teor.~\code{fis:thm:base} em \(m=8\). A identificação
máscara \(=e_k=\) Lei \(k\) está em \code{fis:neuronio}.
\end{proof}

\begin{observacao}[verificação]\label{univ:obs:aranha}
O mesmo enunciado corre sem ramo cujo número de voltas dependa do valor
(\code{fis:neuronio}):
\[
f_k \mathrel{+}= (b\gg k)\wedge 1,\qquad k=0,\dots,7.
\]
Medido em \code{papers/aranha.c} e \code{papers/aranha.asm}
(\texttt{neuronio8}): oito somas fixas, \(\sum_k f_k\) conserva o
popcount. Gram nos \(64\) pares: \code{tests/aranha\_n.c} §AN24.
Isto \textbf{não} institui um corpo de \(256\) elementos como 14º lugar:
é a escrita do espaço que \(\mathcal{U}\) consome.
\end{observacao}

%═══════════════════════════════════════════════════════════════════════════
\section{Operações do corpo, e morfismos entre corpos}\label{univ:sec:modulos}

\begin{definicao}[a cisão que fecha a espinha]\label{univ:def:modulos}
A espinha \textbf{não} é uma cadeia em que um implica o outro, e
\textbf{não} mistura o que \(\mathcal{U}\) \emph{faz} com o que
\(\mathcal{U}\) \emph{classifica}.
\caixa{%
$\begin{array}{rcl}
X,X^{*},\mathsf{Aut}
&:&
\text{espaço, dual, autómato}\\[4pt]
\mathsf{Mor}
&:&
\operatorname{Hom},\ \operatorname{Iso},\ \operatorname{Diff},\ \operatorname{Isom}_{\mu},\ \mathbf{Duo}
\end{array}$}
O primeiro bloco processa. O segundo relaciona. Não se misturam.
\end{definicao}

\begin{definicao}[os quatro contratos de mapa]\label{univ:def:mapas}
Isto \textbf{não} são operações de \(\mathcal{U}\). São os morfismos
que \(\mathcal{U}\) classifica. Postas construções \(A,B\) (ou variedades
\(M,N\)) e \(f\) entre os suportes, \(\mathcal{U}\) pergunta o que \(f\)
preserva --- Def.~\code{fis:def:duomorf} para as duas primeiras linhas
algébricas; as duas últimas quando essa estrutura existe no capítulo.
\begin{center}\small
\begin{tabular}{@{}lll@{}}
\toprule
módulo & o que \(f\) preserva & equação \\
\midrule
\(\operatorname{Hom}(A,B)\)
  & a operação
  & \(f(a\oplus b)=f(a)\oplus f(b)\)
    (e a segunda lei, se houver) \\
\(\operatorname{Iso}(A,B)\)
  & o mesmo, e é bijecção
  & Hom bijetivo; bit \(a=0\) de \code{fis:def:duomorf} \\
\(\operatorname{Diff}(M,N)\)
  & a estrutura diferenciável
  & \(f\) bijecção \(C^{\infty}\) com inversa \(C^{\infty}\) \\
\(\operatorname{Isom}(M,N)\)
  & o volume (canónico de \(\mathcal{U}\))
  & \(\lvert\det Df\rvert=1\), i.e.\ \(\mu(f(E))=\mu(E)\) \\
\(\mathbf{Duo}(A,B)\)
  & troca \(\oplus\) com \(\otimes\)
  & bit \(a=1\): \(f\colon A\xrightarrow{\cong}\mathcal{D}(B)\) \\
\bottomrule
\end{tabular}
\end{center}
O comprimento euclidiano \(g(fx,fy)=g(x,y)\) é \textbf{outro} contrato
(\code{algebrico thm:metrica}: \(\lvert\det\rvert=1\) sem
\(\lVert Tv\rVert=\lVert v\rVert\)). Não é o canónico de \(\mathcal{U}\).
\end{definicao}

\begin{observacao}[homónimo \(\operatorname{Iso}\)]\label{univ:obs:iso}
\(\operatorname{Iso}\) algébrico \textbf{não} é \(\operatorname{Isom}\)
de volume, \textbf{nem} isometria euclidiana, \textbf{nem} a
ultramétrica do índice. I0 aplica-se a este paper: a coincidência de
letras não promove um ao outro (Def.~\ref{univ:def:tres-reguas}).
O catálogo consulta o módulo, não a palavra.
\end{observacao}

\begin{definicao}[realização / transporte]\label{univ:def:realizacao}
\(\mathcal{U}\) tem um operador de \textbf{realização} que decide
\textbf{qual} estrutura \(f\) está a preservar. Não escolhe o módulo
por analogia com outro capítulo. A pergunta é sempre:
\caixa{$A\ \xrightarrow{\;f\;}\ B$ \quad --- \quad que contrato \(f\) satisfaz?}
Se as duas leis se preservam na mesma ordem: Hom (Iso se bijecção).
Se a estrutura diferenciável está escrita e \(f\) a preserva: Diff.
Se \(\lvert\det Df\rvert=1\) (volume): \(\operatorname{Isom}\) canónico de
\(\mathcal{U}\) (\S\ref{univ:sec:metrica}).
Se o bit \(a=1\) e \(f\colon A\xrightarrow{\cong}\mathcal{D}(B)\):
\(\mathbf{Duo}\) --- contrato próprio, não \(\mathcal{C}^{\mathrm{op}}\).
Se \(g(fx,fy)=g(x,y)\): isometria euclidiana --- outro contrato.
Se nenhuma dessas evidências aparece: \texttt{não localizada}, não
«parece isomorfo».
\end{definicao}

%═══════════════════════════════════════════════════════════════════════════
\section{Linguagem consagrada: categorias, e o autómato}\label{univ:sec:lingua}

\begin{observacao}[homónimo «categoria»]\label{univ:obs:categoria}
O catálogo proíbe \emph{categoria nova} como eixo taxonómico
(\code{catalogo}: sem eixo novo, sem etiqueta nova). Isso \textbf{não}
é a teoria das categorias. I0: coincidência de letras. A linguagem consagrada abaixo é \textbf{comparação} com o que
\(\mathcal{U}\) já classifica. \(\mathcal{U}\) \textbf{não} é uma
categoria universal, nem institui \(\mathbf{Cat}\) no bestiário, nem
um 14º lugar.
\end{observacao}

\begin{definicao}[as categorias concretas que \(\mathcal{U}\) classifica]\label{univ:def:cat}
Os morfismos de Def.~\ref{univ:def:mapas} são, na língua consagrada,
morfismos de \textbf{categorias concretas} sobre as construções que a
evidência nomeia --- objectos \(=\) capítulos/construções com suporte;
setas \(=\) contratos de preservação. Não são uma categoria só:
\caixa{%
$\begin{array}{rcl}
\mathbf{Hom}
&:&
\text{preserva }\oplus\text{ (e }\otimes\text{, se houver)}\\[2pt]
\mathbf{Iso}
&:&
\text{subcategoria das bijecções; bit }a=0\\[2pt]
\mathbf{Diff}
&:&
\text{quando }C^{\infty}\text{ está escrita}\\[2pt]
\mathbf{Isom}_{\mu}
&:&
\lvert\det\rvert=1\\[2pt]
\mathbf{Duo}
&:&
f\colon A\xrightarrow{\cong}\mathcal{D}(B);\ \text{bit }a=1
\end{array}$}
O operador de realização (Def.~\ref{univ:def:realizacao}) escolhe
\[
f\in\{\mathrm{Hom},\mathrm{Iso},\mathrm{Diff},\operatorname{Isom}_{\mu},\mathbf{Duo}\}.
\]
Não se pergunta «parece um duomorfismo»; pergunta-se:
\emph{qual é a evidência para o contrato \(\mathbf{Duo}\) neste suporte?}
\[
\boxed{\;\mathbf{Duo}\neq\mathcal{C}^{\mathrm{op}}\;}
\]
A categoria oposta inverte o sentido da seta,
\(A\xrightarrow{f}B\mapsto B\xrightarrow{f^{\mathrm{op}}}A\).
\(\mathbf{Duo}\) pode permanecer no mesmo suporte e trocar as operações:
\[
\mathbf{D}:
\quad
(\oplus,\otimes)
\longrightarrow
(\otimes,\oplus).
\]
Logo \(\mathbf{Duo}\) é uma \textbf{seta com contrato próprio}, não uma
categoria oposta (Def.~\code{fis:def:duomorf}).
\end{definicao}

\begin{observacao}[2-células: não localizada]\label{univ:obs:2cat}
\[
\boxed{\text{2-células: não localizada}}
\]
Não há 2-células escritas neste contrato. Não se promove a realizada.
O percurso
\(\text{Álgebra}\to\text{Topologia}\to\cdots\to\text{Redes}\)
é execução do autómato sobre uma sequência de entradas, não composição
de 2-morfismos. Um funtor \(F\colon\mathbf{Iso}\to\mathbf{Hom}\)
(esquecimento da bijecção) é comparação legítima: Iso não implica Diff.
\end{observacao}

\begin{definicao}[\(\operatorname{INGEST}\) é autómato]\label{univ:def:aut}
O estado de \(\mathsf{Aut}\) é um vector de \(V\): o campo \(G\)
(Def.~\code{fis:def:doisandares}). O Teor.~\code{fis:thm:mult} \textbf{é}
o autómato; o ciclo da aranha lê-o (\code{fis:algoritmo}) --- escrita,
leitura, decisão, fecho; \textbf{sem deliberação}.
\(\operatorname{INGEST}\) instancia o mesmo objecto:
estado \(=\) (coordenada na base \(\{e_k\}\), evidência);
transição \(=\) (pergunta, evidência, resultado, proibição).
Não explorador. \(\texttt{não localizada}\) não vira \(\texttt{N/A}\)
sem prova estrutural.
\end{definicao}

%═══════════════════════════════════════════════════════════════════════════
\section{Métrica canónica: conservação de volume}\label{univ:sec:metrica}

\begin{definicao}[métrica canónica de \(\mathcal{U}\)]\label{univ:def:volume}
A métrica canónica de \(\mathcal{U}\) é a \textbf{conservação de volume}:
\[
\boxed{\;
\lvert\det DT(x)\rvert=1
\quad\Longrightarrow\quad
\mu\bigl(T(E)\bigr)=\mu(E).
\;}
\]
É o factor local da medida (\code{analitico thm:medida-jacobiano}: já o
nomeia como conservação métrica por dualidade neste Corpo Universal).
No reticulado, exacto sem integral:
\([\mathbb{Z}^{2}:T(\mathbb{Z}^{2})]=\lvert\det T\rvert\), donde
\(\lvert\det T\rvert=1\) se e só se \(T\) é sobrejectiva e conserva
\(\mu\) (\code{algebrico thm:metrica}).
\end{definicao}

\begin{observacao}[Lebesgue mede a medida, não o volume]\label{univ:obs:lebesgue}
A métrica de Lebesgue no Algébrico é a diferença simétrica
\[
d(A,B)=\mu(A\,\triangle\,B)
\]
(\code{algebrico thm:metrica}). Toda bijecção é isometria \emph{nessa}
\(d\), e no entanto \(\Phi(A)=A\triangle K\) é isometria que
\textbf{muda} \(\mu\). Donde:
\caixa{conservar a medida é estritamente mais forte do que preservar \(d\)}
\(\lvert\det\rvert=1\) \textbf{não} é condição topológica: é o
jacobiano. Lebesgue no teorema central GLH
(\code{analitico thm:central-energia}) \textbf{casa a ordem}
(\(\sigma\)-aditividade: a ponte não leva valores, leva ordem). Não
institui \(d=\mu(\triangle)\) como canónica de \(\mathcal{U}\), nem
confunde transporte de ordem com conservação de volume.
\end{observacao}

\begin{definicao}[três réguas, e I0 aplica-se]\label{univ:def:tres-reguas}
\[
\boxed{\;
d_{\mathrm{Lebesgue}}
\;\neq\;
\lvert\det\rvert=1
\;\neq\;
d_{\mathrm{ultra}}
\;}
\]
\begin{center}\small
\begin{tabularx}{\textwidth}{@{}l>{\raggedright\arraybackslash}X>{\raggedright\arraybackslash}X@{}}
\toprule
régua & o que mede & cabe em \(\mathcal{U}\) como \\
\midrule
\(d(A,B)=\mu(A\triangle B)\)
  & a medida, não a forma
  & \texttt{não} canónica: isometrias mudam \(\mu\) \\
\(\lvert\det DT\rvert=1\)
  & o volume
  & \textbf{canónica} (Def.~\ref{univ:def:volume}) \\
\(d(a,b)=2^{-\operatorname{prof}(a,b)}\)
  & profundidade de concordância no \textbf{índice}
  & ultramétrica da família / metaindução \\
\(g(fx,fy)=g(x,y)\)
  & comprimento euclidiano
  & contrato à parte; o gato tem volume sem isto \\
\bottomrule
\end{tabularx}
\end{center}
Dois teoremas centrais, e não se fundem: GLH (Hurwitz conta, Gentil
integra, Lebesgue casa a ordem) \textbf{não} é o teorema central da
ultramétrica (a célula é a bola; compor não acumula).
\end{definicao}

\begin{proposicao}[o que a ultramétrica cabe, e o que não cabe]\label{univ:prop:ultra}
A ultramétrica dos endereços
(\code{fis:def:arvore}, \code{fis:lem:ultra})
\[
d(a,c)\le\max\bigl\{d(a,b),d(b,c)\bigr\}
\]
vive no \textbf{índice}, e a bijecção transporta-a
(\code{catalogo sec:completa-global}): \(d_X(x,y):=d_I(R(x),R(y))\).
As bolas particionam (\code{fis:thm:bolas}): a célula \emph{é} a bola.
A travessia não acumula (\code{aranha thm:navega},
\code{aranha cor:global}):
\[
D(\sigma,\rho)\le\max\bigl\{D(\sigma,\tau),D(\tau,\rho)\bigr\}.
\]
Custo de uma cadeia \(=\) o pior passo, qualquer que seja o comprimento.
Isso é a régua da \textbf{metaindução} e de \(\operatorname{INGEST}\):
a regra não menciona o capítulo, e o preço não soma treze auditorias.
\textbf{Não} é a métrica canónica de \(\mathcal{U}\): não vê volume.
Endereçar é geral; conservar \(\mu\) pede o jacobiano.
\end{proposicao}

\begin{observacao}[homónimo «teorema central»]\label{univ:obs:centrais}
I0: «teorema central» em GLH \textbf{não} é «teorema central» da
ultramétrica. O primeiro transporta contar\(\leftrightarrow\)integrar
sem perder ordem. O segundo diz que bolas de raio \(r\) coincidem ou
são disjuntas, e que \(D\) nas representações é ultramétrica. Promover
um ao outro é o mesmo erro que promover uma coincidência de letras.
\end{observacao}

%═══════════════════════════════════════════════════════════════════════════
\section{Indução e metaindução}\label{univ:sec:inducao}

\begin{definicao}[indução]\label{univ:def:inducao}
A indução prova uma propriedade \textbf{dentro} de um suporte:
\[
P(0),\qquad
P(n)\Rightarrow P(n+1)
\quad\Longrightarrow\quad
\forall n\,P(n).
\]
É o passo da torre (\code{algebrico def:inducao}: toda transformação é
um passo de indução; o corpo do passo não menciona o andar).
\end{definicao}

\begin{definicao}[metaindução]\label{univ:def:metainducao}
A metaindução prova uma propriedade \textbf{da família de
construções}:
\caixa{se a própria regra de construção é preservada, a construção inteira é preservada}
Se o passo \(\mathcal{C}(n)\to\mathcal{C}(n+1)\) não menciona \(n\),
então \(\forall n\,\mathcal{C}(n)\) segue da \textbf{forma} do passo
(\code{algebrico thm:meta-inducao}). O resultado não é a tabela dos
andares: é a forma. É por isso que \(\mathcal{U}\) é universal
\textbf{sem} ser um supercorpo: não absorve os corpos; preserva a
regra que os relaciona. A régua dessa família é a ultramétrica dos
endereços (Prop.~\ref{univ:prop:ultra}): o custo é o pior passo, não a
soma.
\end{definicao}

\noindent Donde o catálogo deixa de repetir treze auditorias. A regra
de ingestão é uma só; a metaindução transporta-a a cada capítulo.
\(\operatorname{INGEST}\) é essa realização (\S\ref{univ:sec:motor}).

%═══════════════════════════════════════════════════════════════════════════
\section{\texorpdfstring{Estrela de \(\mathcal{U}\): o duomorfismo}{Estrela do Universal: o duomorfismo}}\label{univ:sec:estrela}

\begin{definicao}[estrela do Universal]\label{univ:def:star}
A estrela de \(\mathcal{U}\) \textbf{não} é a equação
\(x\otimes x=x\oplus 1\) (Teor.~\code{fis:thm:estrela}: condição sobre
um \textbf{elemento}, noutro suporte). A Física já o separa: a estrela
daquele documento \textbf{não} é \(\operatorname{Fix}(\mathcal{D})\)
(\code{fis:def:duomorf}).
A estrela de \(\mathcal{U}\) é o \textbf{duomorfismo}:
\[
\boxed{\;
\operatorname{Star}(\mathcal{U})
=\mathcal{D},
\qquad
f\colon A\ \xrightarrow{\ \cong\ }\ \mathcal{D}(B)=B^{\vee}
\;}
\]
isto é, o bit \(a=1\) de \code{fis:def:duomorf}:
\[
f(x\oplus_A y)=f(x)\otimes_B f(y),
\qquad
f(x\otimes_A y)=f(x)\oplus_B f(y).
\]
\(\mathcal{D}\) não actua sobre os elementos: actua sobre a função
estrutural das duas operações, no mesmo suporte.
\end{definicao}

\begin{observacao}[\(\mathcal{D}\) não é mapa genérico]\label{univ:obs:D}
\(\mathcal{D}\) da Física é \textbf{mais específico} do que qualquer
módulo de \S\ref{univ:sec:modulos}. Não se confunde:
\begin{center}\small
\begin{tabular}{@{}ll@{}}
\toprule
isto & não é isto \\
\midrule
\(\operatorname{Hom}\) & \(\mathcal{D}\) (Hom preserva a ordem das leis; \(\mathcal{D}\) troca-as) \\
\(\operatorname{Iso}\) & \(\mathcal{D}\) (Iso é \(a=0\); duomorfismo é \(a=1\)) \\
\(\operatorname{Diff}\) & \(\mathcal{D}\) (estrutura \(C^{\infty}\), outro suporte) \\
\(\operatorname{Isom}\) & \(\mathcal{D}\) (volume \(\lvert\det\rvert=1\), outro suporte) \\
\(\operatorname{Star}(S)\) de \code{fis:thm:estrela} & \(\operatorname{Star}(\mathcal{U})\) \\
corpo dual \((V,V^{*})\); Poincaré; dual do Rei & \(\mathcal{D}\) do duomorfismo \\
\bottomrule
\end{tabular}
\end{center}
O catálogo já tem esta lista de homónimos. \(\mathcal{U}\) não a
reabre: consome-a em I0.
\end{observacao}

%═══════════════════════════════════════════════════════════════════════════
\section{Vocabulário do resultado}\label{univ:sec:vocab}

\begin{definicao}[cinco valores, e um sexto que não é valor de existência]\label{univ:def:vocab}
Cada slot --- de mapa, de módulo, ou de ingestão --- devolve
\textbf{exactamente um} dos seguintes.
\begin{center}\small
\begin{tabularx}{\textwidth}{@{}l>{\raggedright\arraybackslash}X>{\raggedright\arraybackslash}X@{}}
\toprule
valor & o que significa & o que \emph{não} significa \\
\midrule
\texttt{não localizada}
  & procurou-se; a evidência não apareceu
  & que o objecto não existe \\
\texttt{N/A}
  & demonstrou-se que a estrutura \textbf{não admite} o passo
  & que o agente não achou \\
\texttt{dívida}
  & a hipótese falta (há o que pedir)
  & ausência demonstrada \\
\texttt{candidato}
  & a hipótese está escrita; a realização falta
  & «parece combinar» / «parece isomorfo» \\
\texttt{realizado}
  & evidência matemática preenche o slot
  & analogia \\
\bottomrule
\end{tabularx}
\end{center}
\caixa{não encontrado $\neq$ não existe}
\texttt{não localizada} é valor de \emph{busca}. \texttt{N/A} é valor de
\emph{estrutura}, e só depois de demonstrar que o suporte não admite
aquela etapa. Confundi-los é o erro que promove homónimo a termo técnico.
\end{definicao}

\noindent Primeiro \textbf{procurar}, depois \textbf{demonstrar
ausência}, só então classificar \texttt{N/A}. Se o agente encontra uma
palavra («estrela», \(\mu\), «iso») em \code{redes.tex} ou
noutro sítio, I0 pergunta se é estrutura ou homónimo \textbf{antes} de
pôr o nome na ficha.

%═══════════════════════════════════════════════════════════════════════════
\section{\texorpdfstring{Ingestão: uma realização de \(\mathcal{U}\)}{Ingestão: uma realização do Universal}}\label{univ:sec:contrato}

\begin{definicao}[ingestão é realização, não definição]\label{univ:def:curta}
\(\operatorname{INGEST}\) instancia \(\mathsf{Aut}\) nos capítulos: a
regra não menciona o lugar. Forma curta --- setas verificáveis, não
promessa de existência ---:
\[
\boxed{\;
\operatorname{INGEST}:\quad
S
\to
(\oplus,\otimes)
\to
C
\to
\nu
\to
\operatorname{Star}
\to
\mathcal{D}
\to
\mathrm{fecho}
\to
\mathrm{estado}.
\;}
\]
Aqui \(\operatorname{Star}\) é a de \code{fis:thm:estrela} \emph{no
capítulo}; \(\mathcal{D}\) é o duomorfismo \emph{no capítulo}. Não se
confundem com \(\operatorname{Star}(\mathcal{U})\)
(Def.~\ref{univ:def:star}). A ordem \(\mathrm{I}_{0},\ldots,\mathrm{I}_{11}\)
está em \S\ref{univ:sec:motor}; a evidência fica no catálogo.
\end{definicao}

\begin{observacao}[\(\mathrm{lcm}\) e \(\mathrm{gcd}\) são uma realização]\label{univ:obs:lcm}
A identidade
\[
\otimes=\operatorname{lcm},\qquad
\otimes^{*}=\operatorname{gcd},\qquad
ab=\operatorname{lcm}(a,b)\operatorname{gcd}(a,b)
\]
com \(\nu_N(d)=N/d\) a trocar os papéis (\code{topologico def:somaprod},
\code{topologico thm:reticulo}) é \textbf{uma} realização de
\(\mathrm{I}_{7}\), medida no reticulado \((\mathbb{N}_{>0},\mid)\).
\textbf{Não} é regra de que todo corpo tenha \(\mathrm{lcm}/\mathrm{gcd}\).
A cadeia numérica responde \(\mathcal{D}=\texttt{N/A}\) nesse slot, e
cumpre o contrato. A EM selecciona esta realização para a
\emph{interface}, não para \(V=\mathbb{R}^{n}\) como reticulado.
\end{observacao}

%═══════════════════════════════════════════════════════════════════════════
\section{\texorpdfstring{O motor --- \(\operatorname{INGEST}\)}{O motor --- INGEST}}\label{univ:sec:motor}

\begin{definicao}[o quádruplo de cada estágio]\label{univ:def:quadruplo}
Todo slot \(\mathrm{I}_{k}\) tem quatro campos fixos:
\caixa{(pergunta, evidência, resultado, proibição)}
\begin{itemize}\itemsep2pt
\item \textbf{pergunta} --- o que o motor pergunta, e só isso;
\item \textbf{evidência} --- onde já está escrito (Física, paper, medidor);
\item \textbf{resultado} --- um valor de Def.~\ref{univ:def:vocab}, mais o
conteúdo se \texttt{realizado};
\item \textbf{proibição} --- o que não se pode pôr na ficha mesmo que a
palavra apareça no texto.
\end{itemize}
A proibição é o que impede o homónimo de virar termo técnico: a pergunta
de \(\mathrm{I}_{0}\) corre \textbf{antes} de qualquer promoção.
Na língua do autómato (Def.~\ref{univ:def:aut}): o quádruplo \textbf{é}
a transição; o resultado é um juízo, não uma analogia.
\end{definicao}

\begin{definicao}[\(\operatorname{INGEST}\)]\label{univ:def:ingest}
Dado um capítulo (Parte da Física, paper, ou nome do bestiário),
\[
\operatorname{INGEST}(\textit{capítulo})
\]
percorre \(\mathrm{I}_{0},\ldots,\mathrm{I}_{11}\) \textbf{nesta ordem}
--- máquina de estados, não recomendações --- e devolve ao catálogo o
túplo
\[
(\text{suporte},\
\text{operações},\
\text{construção},\
\text{morfismos},\
\text{autómato},\
\text{estado}).
\]
O catálogo \textbf{não} descobre a matemática: recebe este túplo e
regista.
\caixa{\texttt{fichaingestao} $\to$ calibrador $\to$ \texttt{fichacapitulo}}
Lei complementar só depois (\code{catalogo cat:derivacao}). Não se
recomeça em \(\mathrm{I}_{0}\) se a \texttt{fichaingestao} do capítulo já
existe: lê-se. Não se inventa o objecto que o slot não localizou.
O passo seguinte deste contrato \textbf{não} é outra camada teórica:
é aplicar \(\mathcal{U}\) a \code{fisica.tex}, Parte a Parte.
\end{definicao}

\subsection*{A ordem de \(\mathsf{Aut}\)}
\label{univ:tab:pipeline}

\noindent Não é protocolo à parte: é a ordem em que \(\mathsf{Aut}\) lê
coordenadas. Forma curta: Def.~\ref{univ:def:curta}. Uma linha por passo
--- pergunta, e só isso.

{\small
\begin{enumerate}\itemsep1pt
\item[\(\mathrm{I}_{0}\)] estrutura ou homónimo?
\item[\(\mathrm{I}_{1}\)] onde vive \(S\)?
\item[\(\mathrm{I}_{2}\)] cisão no mesmo \(S\)? (uma só op \(\Rightarrow\) \(\mathrm{I}_{7}\) \texttt{N/A})
\item[\(\mathrm{I}_{3}\)] o que cada lado escreve?
\item[\(\mathrm{I}_{4}\)] as duas contagens somam o rectângulo? (\code{fis:thm:central})
\item[\(\mathrm{I}_{5}\)] transformada, dois polos?
\item[\(\mathrm{I}_{6}\)] Lei~7: ligam sem fundir?
\item[\(\mathrm{I}_{7}\)] que operações, e quem as troca?
\item[\(\mathrm{I}_{8}\)] que involução?
\item[\(\mathrm{I}_{9}\)] \(x^{2}=x+1\) neste \(S\)? (não \(x^{2}=-1\); não \(\operatorname{Star}(\mathcal{U})\))
\item[\(\mathrm{I}_{10}\)] o que fecha, por que mecanismo?
\item[\(\mathrm{I}_{11}\)] par e estado; \(\mathsf{Mor}\) se houver \(f\)
\end{enumerate}}

A evidência, o resultado e a proibição de cada passo estão no catálogo
(\texttt{fichaingestao}). Este paper não os volta a legislar.

\subsection*{Paragem}

\begin{proposicao}[regras de paragem]\label{univ:prop:paragem}
\begin{enumerate}\itemsep2pt
\item \texttt{não localizada} passa ao slot seguinte \textbf{sem} inventar
o objecto e \textbf{sem} promover a \texttt{N/A}.
\item \texttt{N/A} em \(\mathrm{I}_{2}\) (uma só operação) implica
\texttt{N/A} em \(\mathrm{I}_{7}\). Não se pergunta «qual seria a dual».
\item \(\mathrm{I}_{7}\) realizado \textbf{não} obriga \(\mathrm{I}_{10}\)
realizado.
\item Dívida e candidato não se encerram por analogia, nem por «parece
isomorfo».
\item Se a \texttt{fichaingestao} já existe no catálogo, o trabalho é ler,
não reabrir \(\mathrm{I}_{0}\).
\item Classificar \(f\) como Hom / Iso / Diff / Isom / \(\mathcal{D}\)
pede o contrato de Def.~\ref{univ:def:mapas} e Def.~\ref{univ:def:star}.
A palavra no texto não basta.
\end{enumerate}
\end{proposicao}

%═══════════════════════════════════════════════════════════════════════════
\section{O que o agente recebe}\label{univ:sec:agente}

\noindent Não «descubra qual corpo é este». Lê o vector em \(V\), executa
a transição autorizada, pergunta a evidência do contrato de \(f\) se
houver mapa. Saída: o túplo de Def.~\ref{univ:def:ingest}.

\begin{observacao}[autoridade]\label{univ:obs:autoridade}
A Física constrói. \(\mathcal{U}\) opera no espaço e recusa a promoção.
O catálogo regista coordenadas. Os treze correm o mesmo motor
(Def.~\ref{univ:def:metainducao}). EM (\code{catalogo cat:bloco4}) é
preenchimento, não legislação.
\end{observacao}

\end{document}
